% =============================================================================
% 06_CONCLUSION — Concluding Synthesis
% =============================================================================
% Condensed 2026-06-26 from 275 → ~155 words. Detailed implications and
% cornerstone synthesis are treated in Sections 04_thematic and 05_discussion;
% the conclusion restates the central finding and points forward.
% =============================================================================

\section{Conclusion}
\label{sec:conclusion}
\addcontentsline{toc}{section}{Conclusion}

This review has systematically examined the intersection of health AI
evaluation methods and governance structures through a structured content
analysis of the v3 corpus (n=60), surfacing three interdependent
structural deficits --- the \textbf{ecological validity gap} in evaluation
methodology, the \textbf{governance capacity deficit} in health systems,
and the \textbf{pilot trap} --- that jointly account for a substantial share
of documented clinical AI deployment failures. The dominant empirical
cascade runs from ecological validity failure to pilot-trap outcomes, and
none of the four major regulatory frameworks reviewed mandate
ecologically-valid evaluation as a precondition for deployment. Detailed
implications for research, policy, and clinical practice are developed in
Section~\ref{sec:discussion}; the contributions of the four cornerstone
papers that anchor this analysis are synthesized in
Section~\ref{sec:thematic}.

Safe, equitable, and sustainable clinical AI therefore requires an
\textbf{integrated governance-and-evaluation ecosystem} in which technical
validation, organizational capacity-building, regulatory oversight, and
participatory accountability evolve together through implementation --- not
as independent sequential phases. Closing the ecological validity gap is
the most tractable point of intervention, since it conditions both the
evidence on which governance operates and the implementation trajectories
that governance must oversee.
